\documentclass[11pt]{article}
\usepackage{amsmath,amssymb,geometry}
\geometry{margin=1in}
\title{A counterexample to Conjecture 00000008432}
\author{gaochengzhecpu}
\date{}
\begin{document}
\maketitle

The first assertion says that the hypergraph formed by the supports of
minimum-weight codewords has rank $t+1$ when it is a $t$-design. Under the
standard definition, hypergraph rank is the maximum cardinality of an edge.
The binary simplex code of length seven supplies a counterexample: its
minimum-weight supports form a $2$-$(7,4,2)$ design, while their hypergraph
rank is $4$, not $2+1=3$.

\paragraph{The code and all its minimum supports.}
Label the seven coordinates by the nonzero vectors $v\in\mathbb F_2^3$,
and define
\[c(u)=(u\cdot v)_{v\in\mathbb F_2^3\setminus\{0\}},\qquad
 C=\{c(u):u\in\mathbb F_2^3\}.\]
The dot product is over $\mathbb F_2$. The map $c$ is linear and injective:
the coordinates at the three standard basis vectors recover the three
entries of $u$. Thus $C$ is a binary linear code of dimension three.
For $u\ne0$, precisely four vectors $v\in\mathbb F_2^3$ satisfy
$u\cdot v=1$. None of them is zero. All seven nonzero codewords therefore
have weight four; these are exactly the minimum-weight codewords. Their
seven supports are distinct, since binary codewords are determined by
their supports.

For any distinct nonzero coordinates $v,w$, these vectors are linearly
independent over $\mathbb F_2$. The simultaneous equations
$u\cdot v=u\cdot w=1$ have exactly two solutions $u$, both nonzero. Hence
every pair of coordinates lies in exactly two supports, proving the
$2$-$(7,4,2)$ design property. Its strength is not three: if $z=v+w$, the
three equations with right side $1$ are inconsistent, whereas for three
independent coordinate vectors they have exactly one solution. Thus $t=2$
is the actual design strength, and not an arbitrarily chosen weaker
parameter. Since every support has cardinality four, the hypergraph rank
is exactly four. This disproves the rank assertion and therefore the full
conjunction in the conjecture.

\paragraph{Formal verification.}
Lean represents binary triples by $\operatorname{Fin}(8)$ and all length-seven
binary words by $\operatorname{Fin}(128)$. The bit encoding is proved
surjective onto every Boolean coordinate function and every predicate on
the seven coordinates, and injective. The addition and scalar operations
are verified to be the coordinate operations of $\mathbb F_2$.
The encoder is the displayed dot product with all seven nonzero columns.
Its linearity and injectivity, all codeword weights, all minimum supports,
and the absence of repeated blocks are proved.

The pair-count theorem quantifies over every pair of distinct coordinates;
the two explicit triple counts also certify that the strength is not three.
The predicate \texttt{HypergraphRank} requires both a bound for every minimum
support and a support attaining that bound. The final theorem
\texttt{conjecture8432\_counterexample} includes the linear-code and design
hypotheses and proves rank four and failure of rank three. These facts are
kernel-checked finite statements, not an external search certificate.

The proof uses Lean 4.19.0 with \texttt{Std} only and contains no admitted
lemmas, native evaluation proofs, or custom axioms.
\end{document}
